%%%%%%%%%%%%%%%%%%%%%%%%%%%%%%%%%%%%%%%%%
% Journal Article
% LaTeX Template
% Version 2.0 (February 7, 2023)
%
% This template originates from:
% https://www.LaTeXTemplates.com
%
% Author:
% Vel (vel@latextemplates.com)
%
% License:
% CC BY-NC-SA 4.0 (https://creativecommons.org/licenses/by-nc-sa/4.0/)
%
% NOTE: The bibliography needs to be compiled using the biber engine.
%%%%%%%%%%%%%%%%%%%%%%%%%%%%%%%%%%%%%%%%%
% Modified by:
% Pakorn Wangsuekul (joe@ai.iit.tsukuba.ac.jp)
%
% Version 2.1 (July 9, 2023)
% Version 2.2 by Hiroaki Yano (July 11, 2023) 
% Version 2.3 (July 13, 2023)
%%%%%%%%%%%%%%%%%%%%%%%%%%%%%%%%%%%%%%%%%

%----------------------------------------------------------------------------------------
%	PACKAGES AND OTHER DOCUMENT CONFIGURATIONS
%----------------------------------------------------------------------------------------

\documentclass[
	a4paper, % Paper size, use either a4paper or letterpaper
	10pt, % Default font size, can also use 11pt or 12pt, although this is not recommended
	% unnumberedsections, % Uncomment to remove section numbering
	% twoside, % Uncomment to enable two side traditional mode where headers and footers change between odd and even pages
]{LTJournalArticle}

\addbibresource{bibliography.bib} % BibLaTeX bibliography file

\headerleft{Progress Report in PBL I, III/ Research Proposal in PBL II (July/February, 202x)} % Left running head
\headerright{20xxxxxxx Author’s name} % Right running head: Student ID and Name

\footertext{} % Text to appear in the footer, leave this command empty for no footer text

\setcounter{page}{1} % Page number of the first page, set this to a higher number if the article is to be part of a larger publication

%----------------------------------------------------------------------------------------
%	TITLE SECTION
%----------------------------------------------------------------------------------------

% Report title
\title{The Guideline for Writing a Progress Report/ Research Proposal for IDE PBLs} % Article title, max 3 lines recommended

% Author's name
\author{
	Author's Firstname Lastname \hfill \fontsize{9pt}{12pt}\selectfont Bachelor's Program in Interdisciplinary Engineering
}

% Supervisor name and Laboratory name
\date{
        Supervised by: Firstname LASTNAME %\hfill Insert Laboratory Name Here
}

% Abstract (approx. 100 words) and keywords (approx 5 words)
\renewcommand{\maketitlehookd}{%
	\begin{abstract}
            \textbf{Abstract:} This paper explains the guideline for writing a progress report/ research proposal for IDE PBLs. As a standard, it is recommended to write minimum of 2 pages for your report/proposal. The body text of the document should be written in English. The body text should contain the background of the research topic(s), the purpose of the proposed research, the hypothesis for solving the issue(s) mentioned in the research topics, and the verification method of the appropriateness of the hypothesis. In addition, preliminary results and discussion, as well as research implementation plan should preferably be mentioned in the document. Data and references may be included in the appendix section if they cannot be accommodated in the body text.
            
            \textbf{Keywords:} keyword-1, keyword-2, keyword-3, keyword-4, keyword-5
	\end{abstract}
}

%----------------------------------------------------------------------------------------

\begin{document}

\maketitle % Output the title section

%----------------------------------------------------------------------------------------
%	ARTICLE CONTENTS
%----------------------------------------------------------------------------------------

\section{Introduction}

This paper explains the IDE format of progress report and research proposal for PBL I, II, and III in IDE.

\subsection{Pages}
It is recommended to write at least two-page progress report or research proposal.

\subsection{Abstract}
A 100-word abstract should be written in the research proposal. In the progress report, it is not mandatory.

\section{Figures, Tables, and Equations}

\subsection{Figures}
A sample figure is shown below. Using the labels defined in the figure, the figure can be referenced by like this: Figure \ref{fig:sample} shows a sample figure. The figure number is a clickable link, and remains synchronized regardless of other figures inserted before or after it. Full-width figures are shown on the following page.

\begin{figure}[h] % Single column figure
        \centering
	\includegraphics[width=3cm]{Figures/Picture1.png} % Use [width=\linewidth] for full-column figure width
	\caption{Insert figure caption here}
	\label{fig:sample}
\end{figure}

\subsection{Tables}
Similarly, tables can be shown and labeled like figures. Referencing a table using its label: Table \ref{tab:distcounts}.

\begin{table} % Single column table
	\caption{Example single column table.}
	\centering
	\begin{tabular}{l l r}
		\toprule
		\multicolumn{2}{c}{Location} \\
		\cmidrule(r){1-2}
		East Distance & West Distance & Count \\
		\midrule
		100km & 200km & 422 \\
		350km & 1000km & 1833 \\
		600km & 1200km & 890 \\
		\bottomrule
	\end{tabular}
	\label{tab:distcounts}
\end{table}

\subsection{Equations}
Similarly, equations can be shown and labeled like figures and tables. Referencing an equation using its label: Equation \ref{eq:example}.

\begin{equation}
	\cos^3 \theta =\frac{1}{4}\cos\theta+\frac{3}{4}\cos 3\theta
	\label{eq:example}
\end{equation}

%------------------------------------------------

\section{Main Text}
In the body text, the background of the research topic(s), the purpose of the proposed research, the hypothesis for solving the issue(s) mentioned in the research topics, and the verification method of the appropriateness of the hypothesis should be included. In addition, preliminary results, and discussion, as well as research implementation plan should preferably be mentioned.

This line shows how to use a footnote to further explain or cite text\footnote{Example footnote text.}.

\subsection{Unnumbered and Numbered Lists}
This is a bullet point list:

\begin{itemize}
	\item Arcu eros accumsan lorem, at posuere mi diam sit amet tortor
	\item Fusce fermentum, mi sit amet euismod rutrum
        \begin{itemize}
            \item Arcu eros accumsan lorem, at posuere mi diam sit amet tortor
        \end{itemize}
        \item Fusce fermentum, mi sit amet euismod rutrum
\end{itemize}

This is a numbered list:

\begin{enumerate}
	\item Arcu eros accumsan lorem, at posuere mi diam sit amet tortor
	\item Fusce fermentum, mi sit amet euismod rutrum
        \begin{itemize}
            \item Arcu eros accumsan lorem, at posuere mi diam sit amet tortor
        \end{itemize}
        \item Fusce fermentum, mi sit amet euismod rutrum
\end{enumerate}

\subsubsection{Subsubsection Example}

This is an example of the sub-sub-section. In this document, supported levels of sectioning are: section, subsection, and subsubsection.

\section{References}
This template utilizes an IEEE-style citation standard, which includes in-text citations, numbered in square brackets, which refer to the full citation listed in the reference list at the end of the paper. The reference list is organized numerically, not alphabetically.

To cite a source, first, the source's information should be added to the bibliography.bib file as a new entry. Those sources can then be cited here using the cite function. For example: According to \cite{Smith:2024jd}, this hypothesis is incorrect. This statement requires multiple citations \autocite{Smith:2023qr, Smith:2024jd}. This statement contains an in-text citation, for directly referring to a citation like so: \textcite{Smith:2024jd}.

These citations are linked, and clicking them will jump to the references section at the end of the document. For more information regarding citations and references in the LaTeX environment click \href{https://www.overleaf.com/learn/latex/Bibliography_management_in_LaTeX}{\textbf{here}}.

\textbf{To change the bibliography style:} Open the LTJournalArticle.cls file, scroll down to the BIBLIOGRAPHY section. Change the \textit{citestyle} and \textit{bibstyle} options as desired. Information on styles that are available is provided \href{https://www.overleaf.com/learn/latex/Biblatex_bibliography_styles}{\textbf{here}} and \href{https://www.overleaf.com/learn/latex/Biblatex_citation_styles}{\textbf{here}}

\section{Appendix}
Data and references may be included in the appendix section if they cannot be accommodated in the body text. There is no limit to the number of pages in the appendix.

\section{Conclusion}
Submit your progress report for PBL I and III, or research proposal for PBL II in the .pdf format using the Assignment menu of PBL course in manaba.

%----------------------------------------------------------------------------------------
%	 REFERENCES
%----------------------------------------------------------------------------------------

\printbibliography % Output the bibliography

%----------------------------------------------------------------------------------------

\end{document}
